\section{Appendix:  Grammar}

\subsection{Convention}

This specification uses BNF, Perl Compatible Regular Expressions (PCRE), and several trivial augmentations to BNF as
follows:

\begin{itemize}
\item \verb_<>_ surrounds rule names.  When empty, represents the null rule.
\item \verb_""_ means literal
\item \verb_[]_ means optional (0 or 1 occurrence)
\item \verb_()_ means grouping
\item \verb_*_ means 0 or more occurrences
\item \verb_+_ means 1 or more occurrences
\item \verb_//_ means comment
\end{itemize}

\subsection{Token separation}

The compiler uses PCREs to tokenize the input file.  The first capture buffer (\$1) captures a token, and the second
(\$2) captures white-space.  If the expression doesn't match, the compiler tries the next expression.  If all fail, the
compiler reports a syntax error and aborts.  If an expression captures only white-space, matching restarts after the
white-space with the first regular expression.  If an expression captures only a token, it passes the token to the
parser, then again restarts matching after the token with the first regular expression.  No regular expression matches
both a token and white-space.  (These are copied directly from the compiler's source code, with syntactical decoration
removed.  Assume the s flag to anchor \verb_^_ only at the beginning of a string containing newlines.)

\begin{sourcestyle}
\begin{verbatimtab}[\sourcetabsize]
// Single-quoted strings
^('[^']*')()

// Double-quoted strings may contain \"
^("(?:\\.|[^\\"])*")()

// Single-line comments (treated like white-space)
^()(//[^\n]*)

// Multi-line comments (treated like white-space)
^()(/\*.*?\*/)

// White-space
^()([ \t\r\n]+)

// Bare-words and numbers
^([A-Za-z0-9\.\_]+)()

// Single special character tokens
^([\(\)\[\]\{\}\,\@\#\$\;])()

// Multi special character tokens
^([\`\~\!\%\^\&\*\-\=\+\\\|\:\<\>\/\?]+)()
\end{verbatimtab}
\end{sourcestyle}

The above identifies boundaries between tokens, but doesn't identify any meaning for the tokens.  First, consider
individual tokens.  Here the specification again uses PCREs for conciseness, but after this section it is restricted to
BNF.

\begin{sourcestyle}
\begin{verbatimtab}[\sourcetabsize]
<number> ::= <integer> | <float>

<integer> ::= ^[0-9]+$
| ^[A-Z0-9]_[A-Z0-9]+$

<float> ::= ^([1-9A-Z]_)?[0-9A-Z]+
  (\.[0-9A-Z]+)?(e[0-9A-Z]+)?$

\end{verbatimtab}
\end{sourcestyle}

Note that the rules above allow a preceding digit and underscore.  This is the notation for expressing a base besides
decimal.  Any base between 2 and 36 is acceptable.  The base is identified by placing the highest digit for that base
before the underscore.  Here are some examples of \verb_<number>_:

\begin{itemize}
\item A decimal integer:  3735928559
\item Another decimal notation:  9\_3735928559
\item Binary for 3735928559:  1\_11011110101011011011111011101111
\item Hexadecimal for 3735928559:  F\_DEADBEEF
\item Base-36 for 3735928559:  Z\_1PS9WXB
\item A decimal float:  32.5
\item Another decimal notation:  9\_32.5
\item Binary for 32.5:  1\_100000.1
\item Hexadecimal equivalent to 32.5:  F\_20.8
\item Base-36 equivalent to 32.5:  Z\_W.I
\end{itemize}

\begin{sourcestyle}
\begin{verbatimtab}[\sourcetabsize]
<quoted-string> ::= ^'[^']*'$
| ^"(\\.|[^\\"])*"$
\end{verbatimtab}
\end{sourcestyle}

Single-quotes do not allow special encodings, so it is simpler to use them for strings with literal back-slashes, but
it is impossible express single-quotes within a string identified by single-quotes.  Double-quotes allow back-slash
encodings such as \verb_\" \\ \n_ etc.  Encoding is the only difference between the quoting types.  The compiler is
8-bit clean, so a UTF-8 string may appear in quotes with no special support.

\begin{sourcestyle}
\begin{verbatimtab}[\sourcetabsize]
<bare-word> ::= ^[A-Za-z0-9_.]+$

<symbol> ::= "$" <bare-word>
\end{verbatimtab}
\end{sourcestyle}

\verb_<symbol>_ is any entry in the symbol table (any variable or function) defined by source code.  Note that the
leading \$ obviates the need for reserved symbol names.  Language keywords and numbers such as 32, F\_20, if, and asm
are distinct from variables and functions named \$32, \$F\_20, \$if, and \$asm.  This allows future versions of the
language to include new keywords without breaking old code, which is a problem that C suffers from frequently.  While
this means there are no reserved symbol names, names should not begin with an underline because these may be used
internally by the compiler tools.  For example, the symbol ``\_GLOBAL\_OFFSET\_TABLE\_'' is used on many platforms for
dynamic linking.

\begin{sourcestyle}
\begin{verbatimtab}[\sourcetabsize]
<type-def> ::= "@" <bare-word>
\end{verbatimtab}
\end{sourcestyle}

\verb_<type-def>_ is any name defined by source code to refer to a particular data type.  Terminologically, a
\verb_<symbol>_ has an associated \verb_<type-expression>_ (which can be a \verb_<type-def>_), and a \verb_<type-def>_
is a short-hand name for a \verb_<type-expression>_.  Admittedly, both are tracked in the compiler's symbol table, so
this specification can be somewhat confusing distinction if unexplained.  \verb_<symbol>_s are variables and functions.
\verb_<type-def>_s are user-defined types for \verb_<symbol>_s.  As with symbols, names beginning with an underline
should be avoided.

\begin{sourcestyle}
\begin{verbatimtab}[\sourcetabsize]
<jump-label> ::= "#" <bare-word>
\end{verbatimtab}
\end{sourcestyle}

A jump-label is any name defined by source code to refer to an entry point for execution.  These are used for exception
handling and for arbitrary jumps.  (Pit uses the keyword ``jump'' instead of ``goto'' for consistency with assembly
language.)  The scope of a jump-label is limited to the function it is defined in.  These use \# instead of \$.  These
labels may be used in expressions to represent a memory address of executable code, which is useful for in-line
assembly code.  Again, names beginning with an underline should be avoided.  The compiler will almost always need to
define jump labels for flow control constructs; these are always named with an underline prefix so they don't conflict
with user-defined jump-labels.  Unlike GCC, user-defined jump-labels are not renamed because in-line assembler code may
need to reference them.

\subsection{Modules and statements}

\begin{sourcestyle}
\begin{verbatimtab}[\sourcetabsize]
<module> ::= <import-statement>* <statement>*
\end{verbatimtab}
\end{sourcestyle}

An entire source file is a \verb_<module>_.  \verb_<import-statement>_s must come before other \verb_<statement>_s.

\begin{sourcestyle}
\begin{verbatimtab}[\sourcetabsize]
<import-statement> ::= "import" <bare-word> ";"
\end{verbatimtab}
\end{sourcestyle}

\verb_<import-statement>_s simply reference other libraries.  Dots in library names form a hierarchical naming
convention corresponding to namespaces and subdirectories in the library search path.

\begin{sourcestyle}
\begin{verbatimtab}[\sourcetabsize]
<statement> ::= <namespace-statement>
| <asm-statement>
| <conditional-statement>
| <loop-statement>
| <exception-statement>
| <declaration-statement>
| <expression> ";"
| "return" ";"

<statement-block> ::= <statement>
| "{" <statement>* "}"

<namespace-statement> ::= "namespace" ":" <statement-block>
| "namespace" <bare-word> ":" <statement-block>
\end{verbatimtab}
\end{sourcestyle}

A namespace may be anonymous (without a name).  Anonymous namespaces have little effect except that the private symbols
in them may only be accessed by other symbols within the anonymous namespace.  Visibility of public symbols is
unaffected by an anonymous namespace.  A named namespace has the same effect on private symbols, plus the effect that
contained public symbols are all prefixed with the name of the namespace and a dot.

\verb_<asm-statement>_ isn't defined yet.  It will represent syntax for specifying in-line assembler code.

\begin{sourcestyle}
\begin{verbatimtab}[\sourcetabsize]
<conditional-statement> ::= ("if" | "unless") "("
  <expression> ")" <statement-block> ["else"
  <statement-block>]

<loop-statement> ::= "loop" [<bare-word>] ":"
  <statement-block> ["iterate" <statement-block>]
| "again" ";"
| "next" ";"
| "break" ";"

<exception-statement> ::= "try" "#" <bare-word>
  <statement-block>
| "throw" <expression> ";"
| "label" "#" <bare-word> ";"
| "jump" <expression> ";"

<declaration-statement> ::=
  <visibility> <type-expression> <type-def> ";"
| <visibility> <type-expression> <symbol> ";"
| <visibility> <type-expression> <symbol> "prototype" ";"
| <visibility> <type-expression> <symbol> "alias" "=" "$"
  <bare-word> ";"
| <visibility> <type-expression> <symbol> "=" <initializer>

<visibility> ::= "public"
| "private"
\end{verbatimtab}
\end{sourcestyle}

\verb_<declaration-statement>_s define both variables and functions.  Functions are declared using initializer syntax,
since functions are just constant symbols that contain executable code.  Functions are generally a bit of an enigma to
syntax, although they are less so in Pit than in most other languages.  For example, the sizeof operator does not know
how to figure out their size; they are always constant data; the function initializer can only be used in a
declaration-statement, not in an generic expression; etc.

\subsection{Type-expressions}

\begin{sourcestyle}
\begin{verbatimtab}[\sourcetabsize]
<type-expression> ::= <type-expression-bracket>
| <type-expression-bare>

<type-expression-bracket> ::= "[" <type-expression-bare> "]"

<type-expression-bare> ::= <type-qualifier>*
  <type-expression2> <type-expression3>

<type-qualifier> ::= "signed"
| "unsigned"
| "const"
| "unconst"
| "volatile"
| "unvolatile"
| "stateless"
| "unstateless"

<type-expression2> ::= "auto"
| "char"
| "int" ["(" <integer> ")"]
| "ptr" ["(" <type-expression> ")"]
| "float" ["(" <integer> ")"]
| "string" "(" <type-expression> ["," <integer>] ")"
| "struct" "(" (<type-expression> <quoted-string> ",")+ ")"
| "function" "(" (("in" | "out" | "inout") <type-expression>
  "$" <bare-word> ",")* ")"
| "@" <bare-word>  // user-defined type
| "$" <bare-word>  // type of another symbol

<type-expression3> ::= <>
| "->" <type-expression3>
| "<-" <type-expression3>
| "(" ")" <type-expression3>
| "(" <quoted-string> ")" <type-expression3>
| "(" <symbol> ")" <type-expression3>
\end{verbatimtab}
\end{sourcestyle}

Type-expressions are used for declaring and prototyping symbols, type-casting, within other type-expressions, and the
sizeof operator.  Note that two of the alternatives in \verb_<type-expression3>_ allow referencing the type of an
element in a struct or a parameter to a function.  Square brackets are used exclusively to identify a type-expression.
For example, a type-cast converting \$foo to a 128-bit integer would look like this:

\begin{sourcestyle}
\begin{verbatimtab}[\sourcetabsize]
$foo[int(128)]
\end{verbatimtab}
\end{sourcestyle}

In certain cases where the syntax requires a type-expression, the square brackets are required
(\verb_<type-expression-bracket>_ is used for these).  Specifically, they are required by type-casts and the sizeof
operator; they are not required by declaring/prototyping symbols or within other type-expressions.

\subsection{Expressions}

\begin{sourcestyle}
\begin{verbatimtab}[\sourcetabsize]
<expression> ::= <value> [<binary-op> <expression>]

<value> ::= <unary-op> <value>
| <value> <postfix-unary-op>
| "sizeof" <type-expression-bracket>
| "(" <expression> ")"
| "(" "if" <expression> "then" <expression> "else"
  <expression> ")"
| <symbol>
| "error_id"
| "line_number"
| <number>
| <quoted-string>
| <jump-label>
| ["array" | "hash"] "{" (<immediate> ";")* "}"

<unary-op> ::= "-" | "~" | "not"

<binary-op> ::= "=" | "+=" | "-=" | "*=" | "/=" | "/-="
| "/+=" | "/<=" | "/>=" | "%=" | "|=" | "~|=" | "^=" | "~^="
| "&=" | "~&=" | "<<=" | ">>=" | ":" | "+:" | "-:" | "*:"
| "/:" | "/-: | "/+:" | "/<:" | "/>:" | "%:" | "|:" | "~|:"
| "^:" | "~^:" | "&:" | "~&:" | "<<:" | ">>:" | "or" | "nor"
| "xor" | "nxor" | "and" | "nand" | "<" | "<=" | "==" | "!="
| ">=" | ">" | "<=>" | "min" | "max" | "+" | "-" | "*" | "/"
| "/-" | "/+" | "/<" | "/>" | "%" | "|" | "~|" | "^" | "~^"
| "&" | "~&" | "<<" | ">>"

<postfix-unary-op> ::= "->"    // dereference operator
| "<-"                         // reference operator
| <type-expression-bracket>    // typecast
| "(" <expression> ")"         // element of a string
| "(" <quoted-string> ")"      // element of a for struct
| "(" (<expression> ",")* ")"  // function call
\end{verbatimtab}
\end{sourcestyle}

\subsection{Notes}

The grammar is mostly presented in an order that minimizes the number of unexplained rules at each point as read from
top to bottom.  For example, before defining \verb_<expression>_ and referencing many new undefined rules, the grammar
ties off all the other loose ends.

One minor but intriguing invention is the syntax for numbers in bases besides decimal.  This allows easy representation
of any base up to 36, and avoids the confusion caused by the common convention of using a leading zero to designate
octal.

At the moment, for-loop syntax is not included.  There is some debate over the best syntax for this.  Rather than
regretting a bad decision later, this matter is unfinished for the time being.  Meanwhile, the \verb_<loop-statement>_
syntax allows a clause called ``iterate'' that mostly covers what for-loops can do.

As presented, the grammar appears to support function declarations within functions, however the compiler does not
implement this.  Also, there are restrictions as to which alternatives for \verb_<statement>_ are allowed inside a
function and which are allowed outside a function.  This reflects how the proof-of-concept compiler is implemented;
code generation, not parsing, enforces these restrictions.

Some details in this grammar still need attention, however there is a functioning compiler that translates the grammar
as presented here into i386 assembly language with support for auto variables.
